% Working narrative, 2026-09-28. Not a completed submission or a claim of
% validated live/geometry performance. See 03_claims_and_validation.md.
% Citation keys match the accompanying references.bib (arXiv records).
\section{Introduction}

A robot's map can support more than pose estimation: it can provide visual
observations for learned policies and spatial information for navigation
\cite{splat_nav,splatsim}. Gaussian representations make these uses possible
within a renderable scene model, motivating their adoption in robotic mapping
\cite{gs_autonomy}. For a robot acquiring new observations, however, the
relevant object is the map available during acquisition, not only the
reconstruction obtained after processing ends. Its usefulness depends on both
visual fidelity and reliable geometry, which must be developed while the
mapper continues to receive new data.

A central decision is which observations supervise this evolving map.
Keyframes provide an economical basis for tracking and reconstruction, but
the views selected for pose estimation need not exhaust the useful evidence
in the image stream. Coverage deficiencies can remain between selected views,
even when additional images of those regions have already been captured.
Prior systems address this issue through post-acquisition keyframe insertion
or by selecting non-keyframes for additional training
\cite{hi_slam2,online_nvs}. These findings motivate treating mapping
supervision as a separate decision rather than restricting it to the
tracking keyframes.

Expanding supervision introduces a second problem. Previously observed views
can constrain regions that would otherwise receive little attention as the
camera moves on, making historical replay valuable \cite{imap,monogs}.
Yet a growing replay pool also distributes finite updates across an
increasing number of views. Earlier observations accumulate more opportunities
for optimization, while newly admitted views may receive only limited training
before acquisition ends \cite{cartgs}. Thus, access to more observations does
not by itself ensure that their information is incorporated into the map.

Existing work has already developed uncertainty-based view selection,
adaptive keyframe optimization, and map-guided replay under limited compute
\cite{online_nvs,cartgs,eligsir}. We build on this line of work to study a
specific question: how should an RGB--inertial mapper admit additional past
observations and distribute repeated supervision when optimization must stop
with the stream? We distinguish the set of observations retained for possible
use from the small set selected for each update. This distinction allows
historical coverage to remain available without requiring every view to be
optimized at every step.

Our design admits non-keyframe RGB observations according to completed mapping
work and allocates updates among recent keyframes, historical keyframes, and
admitted non-keyframes. The tracking keyframes continue to provide geometric
supervision, while additional RGB observations broaden the supervision
available for appearance fitting. Within the historical pools, a count-based
sampling policy adjusts training opportunities using completed view usage.
We evaluate these choices with causal inputs, held-out views excluded from
training, and no optimizer updates after the final input. Controlled
comparisons separate the effects of additional observations from changes in
the loss and update schedule; rendering quality, geometry, and elapsed time
must be assessed separately.

% Before submission, replace the last evaluation sentences with precisely
% completed results. Current evidence uses frozen causal tracking and matched
% renders, not a concurrent live SLAM validation. Geometry remains unverified.
